%----------------------------------------------------------------------------------------
%   Үйлдвэрлэлийн дадлагын тайлан — МУИС, МТЭС, МКУТ загвар (XeLaTeX)
%   Хөрвүүлэх: latexmk -xelatex main.tex
%----------------------------------------------------------------------------------------
%!TEX TS-program = xelatex
%!TEX encoding = UTF-8 Unicode
\documentclass[11pt,a4paper]{report}

\usepackage{fontspec}
\setlength{\headheight}{15pt}
% Times New Roman нь Linux дээр суугаагүй тул metric-compatible Liberation Serif ашиглав.
% TNR суулгасны дараа доорх хоёр мөрийг сольж болно:
% \setmainfont[Ligatures=TeX]{Times New Roman}
\setmainfont[Ligatures=TeX]{Liberation Serif}
\setsansfont{Liberation Sans}
\setmonofont[Scale=0.9]{Liberation Mono}

\usepackage{geometry}
\usepackage{fancyhdr}
\usepackage{float}
\usepackage{afterpage}
\usepackage{graphicx}
\graphicspath{{./}{figures/}}
\usepackage{amsmath,amssymb,amsbsy}
\usepackage{dcolumn,array}
\usepackage{tocloft}
\usepackage{dics}
\usepackage{nomencl}
\usepackage{upgreek}
\newcommand{\argmin}{\arg\!\min}
\usepackage{mathtools}
\usepackage[hidelinks, unicode]{hyperref}

\usepackage{algorithm}
\usepackage{algpseudocode}

\DeclarePairedDelimiter\abs{\lvert}{\rvert}%
\makeatletter
\usepackage{caption}
\captionsetup[table]{belowskip=0.5pt}
\usepackage{subfiles}

\usepackage{listings}
\renewcommand{\lstlistingname}{Код}
\renewcommand{\lstlistlistingname}{\lstlistingname ын жагсаалт}

\usepackage{longtable}
\usepackage{booktabs}
\usepackage{pdfpages} % тушаал, үнэлгээний скан хуудас оруулахад

% --- Ном зүй: IEEE загвар ---
\usepackage[backend=biber,style=ieee]{biblatex}
\addbibresource{references.bib}

\usepackage{color}
\definecolor{codegreen}{rgb}{0,0.6,0}
\definecolor{codegray}{rgb}{0.5,0.5,0.5}
\definecolor{codepurple}{rgb}{0.58,0,0.82}
\definecolor{backcolour}{rgb}{0.99,0.99,0.99}

\lstdefinestyle{mystyle}{
    basicstyle=\ttfamily\small,
    backgroundcolor=\color{backcolour},
    commentstyle=\color{codegreen},
    keywordstyle=\color{magenta},
    numberstyle=\tiny\color{codegray},
    stringstyle=\color{codepurple},
    breakatwhitespace=false,
    breaklines=true,
    captionpos=b,
    keepspaces=false,
    numbers=left,
    numbersep=10pt,
    showspaces=false,
    showstringspaces=false,
    showtabs=false,
    tabsize=2
}

\lstset{style=mystyle}

\let\oldabs\abs
\def\abs{\@ifstar{\oldabs}{\oldabs*}}
\makeatother
\makenomenclature
\setlength{\emergencystretch}{3em} % шаардлагатай үед мөрийг арай уян болгож багтаана
\begin{document}


%----------------------------------------------------------------------------------------
%   Өөрийн мэдээллээ оруулах хэсэг  (хоосон талбаруудыг бөглөнө үү)
%----------------------------------------------------------------------------------------
\title{\textbf{Үйлдвэрлэлийн дадлага ажлын тайлан}}
\aothura{Ц.~Тэнгис} \sisiIda{22B1NUM6249}
% Дадлага хийсэн байгууллагын нэр:
\organization{................................}
% Удирдагчийн зэрэг цол, овгийн эхний үсэг, нэр (албан тушаал):
\supervisor{................................}
\university{МОНГОЛ УЛСЫН ИХ СУРГУУЛЬ}
\faculty{МЭДЭЭЛЛИЙН ТЕХНОЛОГИ, ЭЛЕКТРОНИКИЙН СУРГУУЛЬ }
\department{МЭДЭЭЛЭЛ, КОМПЬЮТЕРИЙН УХААНЫ ТЭНХИМ}
\cityName{Улаанбаатар}
% Хэвлэгдсэн огноо:
\gradyear{2026 он}


%----------------------------------------------------------------------------------------
%   Нүүр хуудас, гарчиг, жагсаалтууд (main-pre.tex)
%----------------------------------------------------------------------------------------
%----------------------Нүүр хуудастай хамаатай зүйлс----------------------------
\pagenumbering{roman}

\maketitle

\doublespace

% Гарчгийг автоматаар оруулна
\setcounter{tocdepth}{1}
\tableofcontents

% Зургийн жагсаалтыг автоматаар оруулна
\listoffigures

% Хүснэгтийн жагсаалтыг автоматаар оруулна
\listoftables

% % Кодын жагсаалтыг автоматаар оруулна
% \lstlistoflistings

% This puts the word "Page" right justified above everything else.
\newpage
%% \addtocontents{lof}{Зураг~\hfill Хуудас \par}
\newpage
%% \addtocontents{lot}{Хүснэгт~\hfill Хуудас \par}

\renewcommand{\cftlabel}{Зураг}


% Үндсэн бичвэр: нэг мөрийн зайтай (Word-ийн 1.0). Гарчиг, жагсаалтууд
% дээрх \doublespace хэвээр үлдэнэ.
\normalspace
\pagenumbering{arabic}



% --- Хуудасны дугаар голд ---
\pagestyle{fancy}
\fancyhf{}
\fancyfoot[C]{\thepage}
\renewcommand{\headrulewidth}{0pt}
\renewcommand{\footrulewidth}{0pt}
\fancyhfoffset{0pt}


%----------------------------------------------------------------------------------------
%   Удирдамжийн 4-р зүйлд заасан бүрэлдэхүүн:
%   1) Нүүр  2) Удирдагчийн үнэлгээ  3) Захирлын тушаал  4) Төлөвлөгөө  5) Тайлан  6) Хавсралт
%   Гарын үсэгтэй скан хуудсуудыг бэлэн болмогц доорх мөрүүдийг тайлж оруулна:
%----------------------------------------------------------------------------------------
% \includepdf[pages=-]{figures/udirdagchiin-unelgee.pdf}   % Удирдагчийн үнэлгээ, дүгнэлт
% \includepdf[pages=-]{figures/zahirlyn-tushaal.pdf}       % Удирдагчийг томилсон тушаал


%----------------------------------------------------------------------------------------
%   Тайлангийн үндсэн хэсэг
%----------------------------------------------------------------------------------------
\chapter{Дадлагын ажлын төлөвлөгөө}
\subfile{subfiles/plan}

\chapter{Байгууллагын танилцуулга}
\subfile{subfiles/company}

\chapter{Даалгаврын тодорхойлолт ба судалгаа}
\subfile{subfiles/research}

\chapter{Системийн зохиомж}
\subfile{subfiles/design}

\chapter{Хэрэгжүүлэлт}
\subfile{subfiles/implementation}

\chapter{Туршилт ба үр дүн}
\subfile{subfiles/results}

\chapter{Дадлагын хугацаанд эзэмшсэн мэдлэг, ур чадвар}
\subfile{subfiles/skills}

\chapter{Дүгнэлт}
\subfile{subfiles/conclusion}

%----------------------------------------------------------------------------------------
%   Ном зүй (IEEE загвар)
%----------------------------------------------------------------------------------------
% ТҮР: бүх эх сурвалжийг ном зүйд харуулна. Бүлгүүдэд \cite{...} ишлэлүүд
% бичигдсэний дараа доорх мөрийг устгана (зөвхөн ишилсэн нь үлдэнэ).
\nocite{*}
\printbibliography[heading=bibintoc, title={НОМ ЗҮЙ}]

%----------------------------------------------------------------------------------------
%   Хавсралт
%----------------------------------------------------------------------------------------
\appendix
\chapter{Хавсралт}
\subfile{subfiles/appendix}

\end{document}
